\documentclass[10pt,a4paper]{amsart}
\usepackage[margin=19mm]{geometry}
\usepackage{amsmath,amssymb,mathtools}
\usepackage[scr=rsfs,cal=euler]{mathalfa}
\usepackage{xfrac}
\usepackage[unicode,hidelinks,bookmarks=false]{hyperref}
\newcommand{\ie}{i.e.\ }
\newcommand{\cf}{cf.\ }
\newcommand{\resp}{respectively\ }
\newcommand{\Subsection}{\subsection}
\providecommand{\nobreakdash}{\nobreak\hbox{-}}
\setlength{\emergencystretch}{2em}
\multlinegap0pt
\usepackage{amssymb}
\usepackage{xcolor}
\usepackage{mathtools}
\newtheorem{lemma}{Lemma}
\newcommand{\bb}[1]{\mathbb{#1}} \newcommand{\td}[1]{\widetilde{#1}}
\newcommand{\pair}[1]{\langle #1\rangle}
\newcommand{\ppair}[1]{\left\langle #1\right\rangle}
\newcommand{\pr}[1]{\left(#1\right)}
\newcommand{\set}[1]{\left\{#1\right\}}
\newcommand{\tr}{\text{tr}}
\title{Almost isoclinic Lagrangian submanifolds}
\author{Tang-Kai Lee, Mu-Tao Wang}
\date{Source: arXiv:2608.14444v2 (CC BY 4.0)}
\begin{document}
\maketitle

\noindent\emph{Source and licence.} Almost isoclinic Lagrangian submanifolds. Version arXiv:2608.14444v2 (CC BY 4.0), distributed by arXiv under the Creative Commons Attribution 4.0 International licence, http://creativecommons.org/licenses/by/4.0/, as recorded in the arXiv licence metadata for this exact version and shown on the abstract page https://arxiv.org/abs/2608.14444v2. Full licence notice: \texttt{licenses/arxiv-2608.14444-CC-BY-4.0.txt}.\par\smallskip

\noindent\textit{Editorial scope.} Counted unit: the article's lemma lem:T2 (source0.tex lines 2013--2038). The article's complete original proof of it is delivered with it (source0.tex lines 2040--2180, one \texttt{proof} environment of 141 lines). No mathematics is written, repaired or re-derived: the statement and the proof are the article's own lines, in the article's own notation, and no retained line is substituted or re-flowed.  Retained uncounted as the source-local context the proof needs: lines 1979--1990.  Nothing was omitted from any retained span, so the package carries no omission record.  No retained line contains a citation, so the wrapper carries no bibliography and no bibliography entry.  No retained line contains a figure environment, an image-inclusion command or drawing code, so no image is delivered and the package declares no asset dependency.  Editorial formatting only, in addition: an AMS \texttt{amsart} preamble that restates the article's own declarations and macros used by the retained lines, namely the environments \texttt{lemma} on separate counters, together with the standard packages those lines need.  The retained environments print in this document as Lemma 1 in order of appearance, whereas the article prints the same items with the numbering of its own full document.  The complete native source of the article as distributed in the arXiv e-print for this exact version is bundled unchanged as \texttt{sources/207671/source0.tex}.
\begin{lemma}
	[$T_1$'s expansion]
	\label{lem:T1}
	For $A=(a_{jk})_{j,k},$
	\begin{align*}
	T_1
	= 2\,\tr\pr{M (A\wedge A)}
	= 2 \sum_{j<k} \frac 1{1+\lambda_j\lambda_k} 
	\pr{a_{jj}a_{kk} - a_{jk}^2}
	\end{align*}
	and we let $T_1'=T_1$ and $T_1''=0.$
\end{lemma}

\begin{lemma}[$T_2$'s expansion]
	\label{lem:T2}
	For $A=(a_{jk})_{j,k},$
	\begin{align*}T_2
	&= - \operatorname{tr}\big(M (A \wedge S + S \wedge A)\, M (A \wedge S + S \wedge A)\big)\\
	&=- \sum_{j<k}\sum_{p<q}
	\frac 1{1+\lambda_j\lambda_k}
	\frac 1{1+\lambda_p\lambda_q}
	\pr{a_{jp}\lambda_k\delta_{kq} - a_{jq}\lambda_k\delta_{kp}
		+ \lambda_j\delta_{jp}a_{kq}
		- \lambda_j\delta_{jq}a_{kp}}^2\\
	&=T_2'+T_2'',
	\end{align*}
	where 
	\begin{align*}T_2'&=- \sum_{j<k}
	\frac 1{(1+\lambda_j\lambda_k)^2} 
	\pr{\lambda_k a_{jj}
		+ \lambda_j a_{kk}
	}^2\text{ and}\\
	T_2''
	&=-\sum_{j<k}\sum_{\substack{m\neq j,k}}
	\frac{2\lambda_m^2}
	{(1+\lambda_j\lambda_m)(1+\lambda_k\lambda_m)}
	\, a_{jk}^2.
	\end{align*}
\end{lemma}

\begin{proof}
	As in Lemma~\ref{lem:T1}, we work with an orthonormal basis $\set{v_j}_{j=1,\cdots,n}$ which diagonalizes $S={\rm diag}(\lambda_1,\cdots,\lambda_n)$ and write $A=(a_{jk})_{j,k}$ in the same orthonormal basis.
	Since $v_j\wedge v_k$'s form a basis of $\wedge^2\bb R^n$ for $j<k,$
	the $(jk,pq)$-entry of $A\wedge S+ S\wedge A$ is
	\begin{align*}
	C_{jkpq}
	:= \pair{(A\wedge S+ S\wedge A)(v_j\wedge v_k), v_p\wedge v_q}
	& = \ppair{
		\sum_c a_{jc}v_c\wedge \lambda_kv_k
		+ \lambda_jv_j\wedge \sum_d a_{kd}v_d,
		v_p\wedge v_q
	}\\
	& = a_{jp}\lambda_k\delta_{kq} - a_{jq}\lambda_k\delta_{kp}
	+ \lambda_j\delta_{jp}a_{kq}
	- \lambda_j\delta_{jq}a_{kp},
	\end{align*}
	and hence it is symmetric in $(jk,pq)$ in the sense that
	\begin{align*}
	C_{pqjk}
	& = a_{pj}\lambda_q\delta_{qk} - a_{pk}\lambda_q\delta_{qj}
	+ \lambda_p\delta_{pj}a_{qk}
	- \lambda_p\delta_{pk}a_{qj}
	= C_{jkpq}.
	\end{align*}
	Thus, if we let $m_{jk}:=1/(1+\lambda_j\lambda_k),$ then
	\begin{align*}
	\tr\pr{M (A \wedge S + S \wedge A) \, M (A \wedge S + S \wedge A) }
	= & \sum_{j<k, p<q}
	m_{jk} C_{jkpq} m_{pq} C_{pqjk}\\
	= & \sum_{j<k, p<q}
	m_{jk} m_{pq}
	\pr{a_{jp}\lambda_k\delta_{kq} - a_{jq}\lambda_k\delta_{kp}
		+ \lambda_j\delta_{jp}a_{kq}
		- \lambda_j\delta_{jq}a_{kp}}^2.
	\end{align*}
	This proves the first identity of the lemma.
	
	Next, we split $T_2$ into two terms $T_2'$ and $T_2''$ by
	\begin{align*}
	T_2'
	&:=- \sum_{\substack{j<k, p<q\\(j,k)= (p,q)}}
	\frac 1{1+\lambda_j\lambda_k}
	\frac 1{1+\lambda_p\lambda_q}
	\pr{a_{jp}\lambda_k\delta_{kq} - a_{jq}\lambda_k\delta_{kp}
		+ \lambda_j\delta_{jp}a_{kq}
		- \lambda_j\delta_{jq}a_{kp}}^2
	\text{ and}\\
	T_2''
	&:= -\sum_{\substack{j<k, p<q\\(j,k)\neq (p,q)}}
	\frac 1{1+\lambda_j\lambda_k}
	\frac 1{1+\lambda_p\lambda_q}
	\pr{a_{jp}\lambda_k\delta_{kq} - a_{jq}\lambda_k\delta_{kp}
		+ \lambda_j\delta_{jp}a_{kq}
		- \lambda_j\delta_{jq}a_{kp}}^2.
	\end{align*}
	The expression of $T_2'$ claimed in the lemma then follows directly by $\delta_{jj}=\delta_{kk}=1$ and $\delta_{jk}=0.$
	
	
	$T_2''$ can be rewritten as
	\begin{align*}
	T_2''
	=
	-\sum_{\substack{\ell<m,\ p<q\\(\ell,m)\neq (p,q)}}
	\frac 1{1+\lambda_\ell\lambda_m}
	\frac 1{1+\lambda_p\lambda_q}
	\pr{
		a_{\ell p}\lambda_m\delta_{mq}
		- a_{\ell q}\lambda_m\delta_{mp}
		+ \lambda_\ell\delta_{\ell p}a_{mq}
		- \lambda_\ell\delta_{\ell q}a_{mp}
	}^2.	
	\end{align*}
	Note that since $(\ell,m)\neq (p,q),$ the pairs will contribute only when $\set{\ell,m}\cap \set{p,q}$ has exactly one element.
	Thus, when we fix two numbers $j<k,$ we have the following four cases.
	We let $r$ be the common index, which will be different from $j$ or $k.$
	
	When $\ell=p=r,$ the third term in the parenthesis is the only non-zero term and we get the term containing $a_{jk}^2$ when $\set{m,q}=\set{j,k}.$
	This then leads to
	\begin{align}\label{T2-1}
	\sum_{\substack{r\neq j,k\\r<j,r<k}}
	\frac 1{1+\lambda_r\lambda_j}
	\frac 1{1+\lambda_r\lambda_k}
	\lambda_r^2 a_{jk}^2
	+ \sum_{\substack{r\neq j,k\\r<k,r<j}}
	\frac 1{1+\lambda_r\lambda_k}
	\frac 1{1+\lambda_r\lambda_j}
	\lambda_r^2 a_{kj}^2.
	\end{align}
	When $\ell=q=r,$ the fourth term in the parenthesis is the only non-zero term and we get the term containing $a_{jk}^2$ when $\set{m,p}=\set{j,k}.$
	This then leads to
	\begin{align}\label{T2-2}
	\sum_{\substack{r\neq j,k\\k<r<j}}
	\frac 1{1+\lambda_r\lambda_j}
	\frac 1{1+\lambda_k\lambda_r}
	\lambda_r^2 a_{jk}^2
	+ \sum_{\substack{r\neq j,k\\j<r<k}}
	\frac 1{1+\lambda_r\lambda_k}
	\frac 1{1+\lambda_j\lambda_r}
	\lambda_r^2 a_{kj}^2.
	\end{align}
	When $m=p=r,$ the third term in the parenthesis is the only non-zero term and we get the term containing $a_{jk}^2$ when $\set{\ell,q}=\set{j,k}.$
	This then leads to
	\begin{align}\label{T2-3}
	\sum_{\substack{r\neq j,k\\j<r<k}}
	\frac 1{1+\lambda_r\lambda_j}
	\frac 1{1+\lambda_r\lambda_k}
	\lambda_r^2 a_{jk}^2
	+ \sum_{\substack{r\neq j,k\\k<r<j}}
	\frac 1{1+\lambda_r\lambda_k}
	\frac 1{1+\lambda_r\lambda_j}
	\lambda_r^2 a_{kj}^2.
	\end{align}
	When $m=q=r,$ the third term in the parenthesis is the only non-zero term and we get the term containing $a_{jk}^2$ when $\set{\ell,p}=\set{j,k}.$
	This then leads to
	\begin{align}\label{T2-4}
	\sum_{\substack{r\neq j,k\\r>j,r>k}}
	\frac 1{1+\lambda_r\lambda_j}
	\frac 1{1+\lambda_r\lambda_k}
	\lambda_r^2 a_{jk}^2
	+ \sum_{\substack{r\neq j,k\\r>k,r>j}}
	\frac 1{1+\lambda_r\lambda_k}
	\frac 1{1+\lambda_r\lambda_j}
	\lambda_r^2 a_{kj}^2.
	\end{align}
	Combining these four terms \eqref{T2-1}, \eqref{T2-2}, \eqref{T2-3}, and \eqref{T2-4}, we get exactly two summations over $r\neq j,k,$ so the terms involving $a_{jk}^2$ in $T_2''$ are 
	\begin{align*}
	\sum_{r\neq j,k}
	\frac 1{1+\lambda_r\lambda_j}
	\frac 1{1+\lambda_r\lambda_k}
	\lambda_r a_{jk}^2
	+ \sum_{r\neq j,k}
	\frac 1{1+\lambda_r\lambda_k}
	\frac 1{1+\lambda_r\lambda_j}
	\lambda_r a_{kj}^2
	= \sum_{r\neq j,k}
	\frac{2\lambda_r^2}
	{(1+\lambda_j\lambda_r)(1+\lambda_k\lambda_r)}
	\, a_{jk}^2.
	\end{align*}
	This completes the proof of the lemma.
\end{proof}

\end{document}
